\documentclass[10pt,a4paper]{amsart}
\usepackage[margin=19mm]{geometry}
\usepackage{amsmath,amssymb,mathtools}
\usepackage[scr=rsfs,cal=euler]{mathalfa}
\usepackage{xfrac}
\usepackage[unicode,hidelinks,bookmarks=false]{hyperref}
\newcommand{\ie}{i.e.\ }
\newcommand{\cf}{cf.\ }
\newcommand{\resp}{respectively\ }
\newcommand{\Subsection}{\subsection}
\providecommand{\nobreakdash}{\nobreak\hbox{-}}
\setlength{\emergencystretch}{2em}
\multlinegap0pt
\usepackage{mathtools}
\usepackage{amssymb}
\usepackage{xcolor}
\usepackage{tikz}
\newtheorem{theorem}{Theorem}
\newcommand{\diver}{{{\mathrm{div}}}}
\title{Dirichlet--Neumann duality for the Basic Spectrum of Riemannian Submersions: A Supersymmetric Perspective}
\author{Vicent Gimeno i Garcia, Paulo Henryque da Costa Silva}
\date{Source: arXiv:2606.12009v1 (CC BY 4.0)}
\begin{document}
\maketitle

\noindent\emph{Source and licence.} Dirichlet--Neumann duality for the Basic Spectrum of Riemannian Submersions: A Supersymmetric Perspective. Version arXiv:2606.12009v1 (CC BY 4.0), distributed by arXiv under the Creative Commons Attribution 4.0 International licence, http://creativecommons.org/licenses/by/4.0/, as recorded in the arXiv licence metadata for this exact version and shown on the abstract page https://arxiv.org/abs/2606.12009v1. Full licence notice: \texttt{licenses/arxiv-2606.12009-CC-BY-4.0.txt}.\par\smallskip

\noindent\textit{Editorial scope.} Counted unit: the article's theorem DirichletNeumannProblem (source0.tex lines 504--540). The article's complete original proof of it is delivered with it (source0.tex lines 542--587, one \texttt{proof} environment of 46 lines). No mathematics is written, repaired or re-derived: the statement and the proof are the article's own lines, in the article's own notation, and no retained line is substituted or re-flowed.  No further source-local context is needed: every cross-reference of every retained line resolves inside the two delivered spans.  Nothing was omitted from any retained span, so the package carries no omission record.  No retained line contains a citation, so the wrapper carries no bibliography and no bibliography entry.  No retained line contains a figure environment, an image-inclusion command or drawing code, so no image is delivered and the package declares no asset dependency.  Editorial formatting only, in addition: an AMS \texttt{amsart} preamble that restates the article's own declarations and macros used by the retained lines, namely the environments \texttt{theorem} on separate counters, together with the standard packages those lines need.  The retained environments print in this document as Theorem 1 in order of appearance, whereas the article prints the same items with the numbering of its own full document.  The complete native source of the article as distributed in the arXiv e-print for this exact version is bundled unchanged as \texttt{sources/202596/source0.tex}.
\begin{theorem}\label{DirichletNeumannProblem}
Let $\Omega\subset B$ be a precompact domain with smooth boundary $\partial \Omega$. Consider the Dirichlet and divergence-free eigenvalue problems for the inversely weighted partners $\mathscr{L}_S^{+}$ and $\mathscr{L}_S^{-}$:
\begin{subequations}\label{DN}
\begin{equation}\label{DN-a}
\left\{
\begin{aligned}
\mathscr{L}_S^{+}\varphi &= \lambda \varphi \quad &&\text{in } \Omega
,\\
\varphi &= 0 \quad &&\text{on } \partial \Omega,
\end{aligned}
\right.
\end{equation}

\begin{equation}\label{DN-b}
\left\{
\begin{aligned}
\mathscr{L}_S^{-}X &= \lambda X \quad &&\text{in } \Omega,\\
\diver(X) &= 0 \quad &&\text{on } \partial \Omega.
\end{aligned}
\right.
\end{equation}
\end{subequations}
Then the following statements hold:
\begin{enumerate}
 \item If $\lambda \neq 0$ and $\varphi \in C^{\infty}(B)$ is a nontrivial solution of the Dirichlet eigenvalue problem for $\mathscr L_S^{+}$, then
 \[
 X \coloneqq \mathcal A_S \varphi = S \nabla \varphi 
 \]
is a nontrivial solution of the divergence-free boundary problem associated with $\mathscr L_S^{-}$ in \eqref{DN-b}, with the same eigenvalue $\lambda$.
 \item If $\lambda \neq 0$ and $X \in \mathfrak{X}(M)$ is a nontrivial solution of the divergence-free boundary problem for $\mathscr L_S^{-}$, then
 \[
 \varphi:=\mathcal{A}^\dag_SX=\frac{1}{S}\diver(X)
 \]
 is a nontrivial solution of the Dirichlet problem associated with $\mathscr L_S^{+}$ in \eqref{DN-a}, with the same eigenvalue $\lambda$.
\end{enumerate}
Consequently, the nonzero Dirichlet spectrum of $\mathscr L_S^{+}$ coincides with the nonzero divergence-free boundary spectrum of $\mathscr L_S^{-}$, counting multiplicities. 
\end{theorem}

\begin{proof}
Assume first that $\varphi$ is a nontrivial solution of \eqref{DN-a} with $\lambda \neq 0$, and define $X \coloneq \mathcal A_S\varphi$. By the intertwining relation
\[
\mathcal A_S\mathscr L_S^{+} = \mathscr L_S^{-}\mathcal A_S,
\]
this yields
\[
\mathscr L_S^{-}X
= \mathscr L_S^{-}(\mathcal A_S\varphi)
= \mathcal A_S(\mathscr L_S^{+}\varphi)
= \lambda \mathcal A_S\varphi
= \lambda X.
\]
It remains to determine the effect on the boundary condition. Since $\diver(X)= \diver(S\nabla\varphi)$ and
\[
\mathscr{L}_S^{+} \varphi \coloneqq \dfrac{1}{S}\diver\left({S}\nabla\varphi \right),
\]
it follows that
\[
\diver(X) = S\mathscr L_S^{+}\varphi = \lambda S\varphi.
\]
Since $S$ is continuous in $\Omega$ and $\varphi|_{\partial \Omega}=0$, we conclude that $\diver(X)=0$ on $\partial \Omega$. Hence, $X$ satisfies the divergence-free boundary problem \eqref{DN-b} for $\mathscr L_S^{-}$. Moreover, if $X \equiv 0$, then $S \nabla\varphi\equiv 0$, and since $S>0$, therefore $\nabla \varphi\equiv 0$. Thus, $\varphi$ is constant in each connected component. Hence, by the Dirichlet boundary condition $\varphi = 0$, a contradiction. Therefore $X \not\equiv 0$.

Conversely, now assume that $X$ is a nontrivial solution of the divergence-free boundary problem \eqref{DN-b} with $\lambda \neq 0$, and define $\varphi \coloneq \mathcal{A}^\dag_SX$. Using again the intertwining relation
\[
\mathcal{A}^\dag_S\mathscr{L}_S^{-} = \mathscr{L}_S^{+}\mathcal{A}^\dag_S,
\]
we conclude
\[
\mathscr L_S^{+}\varphi = \mathscr L_S^{+}(\mathcal A^\dag_SX) = \mathcal A^\dag_S(\mathscr L_S^{-}X) =\lambda \mathcal{A}^\dag_SX = \lambda\varphi.
\]

Furthermore, $\varphi|_{\partial \Omega}= \mathcal{A}^\dag_S (X)|_{\partial \Omega} = (1/S) \diver(X)|_{\partial \Omega} = 0$, so $\varphi$ satisfies the Dirichlet problem \eqref{DN-a}. Finally, if $\varphi = 0$, then $\diver(X)=0$. Hence,
\[
0= \nabla\left( \diver(X) \right) - \diver(X) \nabla \log S =\mathscr L_S^{-}X= \lambda X.
\]
Since $\lambda \neq 0$, this implies that $X \equiv 0$, again a contradiction. Therefore $\varphi \neq 0$.

It follows that the maps
\[
X\mapsto \mathcal A_S^\dag X,
\qquad
\psi\mapsto \mathcal A_S\psi,
\]
define mutually inverse correspondences between the eigenspaces associated with each nonzero eigenvalue $\lambda$. This proves the claimed spectral correspondence.
\end{proof}

\end{document}
